\documentclass[11pt]{article}
% Shared preamble for all daily exercises
% Update this file to change settings (padding, parindent, etc.) for all exercises

\usepackage[margin=0.75in]{geometry}
\usepackage{amsmath}
\usepackage{amssymb}

% Custom settings
\setlength{\parindent}{0pt}
\setlength{\parskip}{0.2cm}

\title{Daily Exercise}
\author{Dhruman Gupta}
\date{4th April}

\begin{document}

% Common header for daily exercises
% This file can be included in other LaTeX documents

\maketitle

\noindent
\textbf{Course:} CS-3330, Theory of Computation

\vspace{0.2cm}

\noindent
\textbf{Question:}\noindent 
The NFA intersection problem is: given two NFA $A_1$ and $A_2$, is $L(A_1) \cap L(A_2) \neq \emptyset$?

Prove that this decision problem is in $P$.

\textbf{Answer:}

Let
\[
A_1 = (Q_1, \Sigma, \Delta_1, S_1, F_1), \qquad A_2 = (Q_2, \Sigma, \Delta_2, S_2, F_2)
\]
be the two NFAs.

Construct the product NFA
\[
A = (Q_1 \times Q_2, \Sigma, \Delta, S_1 \times S_2, F_1 \times F_2),
\]
where
\[
(q_1', q_2') \in \Delta((q_1, q_2), a)
\iff
q_1' \in \Delta_1(q_1, a) \text{ and } q_2' \in \Delta_2(q_2, a).
\]

This is the standard product construction, and it is clear that
\[
L(A) = L(A_1) \cap L(A_2).
\]

So it remains to decide whether $L(A) \neq \emptyset$.

Now, emptiness of an NFA can be checked in polynomial time. The idea is to see whether some accepting state is reachable from some start state. This can be done by a dynamic programming construction: for each state $q$, keep track of whether $q$ is reachable from some start state. Initially, mark all states in $S_1 \times S_2$ as reachable. Then repeatedly propagate reachability along transitions until no new states are added. At the end, if some state in $F_1 \times F_2$ is reachable, then $L(A) \neq \emptyset$; otherwise, $L(A) = \emptyset$.

The number of states in $A$ is $|Q_1||Q_2|$, and the number of transitions is at most polynomial in the sizes of $A_1$ and $A_2$. The product construction is therefore polynomial, and the dynamic programming reachability check is also polynomial.

Hence, we can decide whether
\[
L(A_1) \cap L(A_2) \neq \emptyset
\]
in polynomial time.

Therefore the NFA intersection problem is in $P$.

\end{document}
